
\subsection*{Note pour la page 1.}

La théorie des opérations $\Lambda^i$ dans $D^-(X)$, et des opérations $\lambda^i$ dans $K^\bullet(X)$, a été écrite par P. Deligne en 1968. Il prouve que $K^\bullet(X)$ devient ainsi un pré-$\lambda$-anneau, et lorsque $\underline{O}_X$ est un faisceau d'Algèbres sur un corps, que $K^\bullet(X)$ est même un $\lambda$-anneau. Il est plausible qu'un raffinement de la méthode de Deligne permette de se débarrasser de l'hypothèse restrictive sur $\underline{O}_X$.

\subsection*{Notes pour la page 11}

\noindent(*) (Avril 1969) D. MUMFORD nous signale une démonstration de (4.6) dans l'\uline{anneau de Chow} (sans négliger la torsion), qu'il a obtenue dans un séminaire à Harvard en 1959 ! Le lecteur trouvera la démonstration de ce beau résultat dans SGA 5 VII 9. Ce résultat contient évidemment le résultat cohomologique correspondant, signalé dans la note (**). Cependant, le résultat cohomologique est prouvé par J. P. JOUANOLOU sous des conditions plus générales (sans hypothèse quasi-projective, et dans le cas lisse relatif sur une base quelconque).

\medskip

\noindent(**) (janvier 1969) J. P. JOUANOLOU vient de prouver ces formules pour tout $y$ sans négliger la torsion, cf. SGA 5 VII \S{} 4.9. DELIGNE et MANIN 17 avaient auparavant remarqué indépendamment, et par des arguments (faciles) de nature assez différente, qu'il suffisait d'établir des formules pour $y$ ``algébrique'', et même pour $y=1$ dans 16, pour conclure la même relation dans le cas général.



% Final-index helper macros
\providecommand{\Qcoh}{\operatorname{Qcoh}}
\providecommand{\Coh}{\operatorname{Coh}}
\providecommand{\Parf}{\operatorname{Parf}}
\providecommand{\Lib}{\operatorname{Lib}}
\providecommand{\Loclib}{\operatorname{Loclib}}
\providecommand{\Proj}{\operatorname{Proj}}
\providecommand{\Mod}{\operatorname{Mod}}
\renewcommand{\Ob}{\operatorname{Ob}}
\providecommand{\Hom}{\operatorname{Hom}}
\providecommand{\toramp}{\operatorname{tor.amp}}
\providecommand{\parfamp}{\operatorname{par.amp}}
\providecommand{\coh}{\operatorname{coh}}
\providecommand{\qcoh}{\operatorname{qcoh}}
\providecommand{\parf}{\operatorname{parf}}
\providecommand{\torf}{\operatorname{torf}}
\providecommand{\amp}{\operatorname{amp}}
\providecommand{\num}{\operatorname{num}}
\providecommand{\naif}{\operatorname{naif}}
\providecommand{\alg}{\operatorname{alg}}
\providecommand{\Pic}{\operatorname{Pic}}
\providecommand{\Gr}{\operatorname{Gr}}
\providecommand{\Fil}{\operatorname{Fil}}
\providecommand{\Ch}{\operatorname{Ch}}
\providecommand{\Todd}{\operatorname{Todd}}
\providecommand{\Diff}{\operatorname{Diff}}
\providecommand{\Dcat}{\mathrm D}
\providecommand{\Tr}{\operatorname{Tr}}
\providecommand{\Z}{\mathbb Z}
\providecommand{\Q}{\mathbb Q}

\newpage
